%%% FIELD statement-local-one-step
For every \(\mathbf P=(P_n)\in
\mathcal B_{\epsilon,M,c_0,v_0,v_1}\), the explicit cross-fitted saturated-cell
one-step estimator in \(def:saturated\text{-}one\text{-}step\) satisfies
\[
 \widehat\tau^{\mathrm{OS}}-\tau(P_n)
 =\frac1n\sum_{i=1}^n\phi_{P_n}(O_i)
   +o_{Q_{P_n}^{(n)}}(n^{-1/2}),
\]
\[
 \frac{\widehat V^{\mathrm{OS}}}{V_{\mathrm{eff}}(P_n)}
 \longrightarrow1
 \quad\text{in }Q_{P_n}^{(n)}\text{-probability},
 \qquad
 \frac{\sqrt n\{\widehat\tau^{\mathrm{OS}}-\tau(P_n)\}}
      {\sqrt{V_{\mathrm{eff}}(P_n)}}
 \Rightarrow N(0,1).
\]

%%% FIELD statement-order-one-clt
Let \(\mathbf P=(P_n)\) be a triangular sequence for which, in row \(n\),
\(O_1,\ldots,O_n\) are independent with common law \((P_n)_O\).  If
\(E_{P_n}[\phi_{P_n}(O)]=0\),
\(0<V_{\mathrm{eff}}(P_n)=E_{P_n}[\phi_{P_n}(O)^2]\), and
\[
 \sup_n\frac{E_{P_n}[|\phi_{P_n}(O)|^4]}
                 {V_{\mathrm{eff}}(P_n)^2}<\infty,
\]
then
\[
 \frac{\sum_{i=1}^n\phi_{P_n}(O_i)}
 {\left\{\sum_{i=1}^n
 [\phi_{P_n}(O_i)-n^{-1}\sum_{j=1}^n\phi_{P_n}(O_j)]^2\right\}^{1/2}}
 \Rightarrow N(0,1).
\]

%%% FIELD proof-local-one-step
Fix \(\mathbf P=(P_n)\in
\mathcal B_{\epsilon,M,c_0,v_0,v_1}\), and suppress the row subscript on
\(p_k,\pi_k,\mu_{ak}\).  All constants below may depend on
\(\epsilon,c_0,K_0\), but not on \(n,d_n\), or the chosen member of the
triangular class.  Membership in \(def:balanced\text{-}local\text{-}class\)
gives, separately, iid sampling, fixed overlap, the cell-mean envelope, the
conditional fourth-central-moment envelope, the lower cell-mass bound,
\(d_n=o(\sqrt n)\), and
\[
 v_0M^2\le V_{\mathrm{eff}}(P_n)\le v_1M^2.
\tag{1}
\]

First derive the nuisance rates rather than assuming them.  For a validation
fold \(I_s\), write \(m_s=n-|I_s|\) for its training size.  Since \(K_0\) is
fixed and the fold sizes differ by at most one, \(m_s\ge c_Kn\) for a constant
\(c_K>0\).  Let \(N_{ak,-s}\) and \(N_{k,-s}\) be the corresponding training
arm-cell and cell counts.  Cauchy--Schwarz and
\(ass:conditional\text{-}fourth\text{-}moment\) give
\[
 E_{P_n}[(Y-\mu_{ak})^2\mid X=k,A=a]
 \le
 \{E_{P_n}[(Y-\mu_{ak})^4\mid X=k,A=a]\}^{1/2}
 \le M^2.
\tag{2}
\]
If \(N\sim\operatorname{Binomial}(m,q)\), direct integration of the binomial
probability generating function gives
\[
 E[(N+1)^{-1}]=\frac{1-(1-q)^{m+1}}{(m+1)q},
 \qquad
 N^{-1}\mathbf1\{N>0\}\le2(N+1)^{-1}.
\tag{3}
\]
Conditionally on \(N_{ak,-s}=r>0\), the fitted cell mean is the average of
\(r\) observations from the conditional arm-cell law.  By (2), its conditional
mean-squared error is at most \(M^2/r\).  When \(r=0\), the zero fallback in
\(def:saturated\text{-}one\text{-}step\) has squared error at most \(M^2\) by
\(ass:cell\text{-}mean\text{-}envelope\).  With
\(q_{ak}=P_n(X=k,A=a)\ge\epsilon p_k\), equations (2)--(3) therefore imply
\[
\begin{aligned}
 E\!\left[\sum_{k=1}^{d_n}p_k
  (\widehat\mu_{ak}^{(-s)}-\mu_{ak})^2\right]
 &\le M^2\sum_{k=1}^{d_n}p_k
 \left\{e^{-m_sq_{ak}}+\frac{2}{(m_s+1)q_{ak}}\right\}\\
 &\le C_\epsilon M^2
 \left\{e^{-c_\epsilon m_s/d_n}+\frac{d_n}{m_s}\right\}.
\end{aligned}
\tag{4}
\]
The last inequality uses \(p_k\ge c_0/d_n\) and
\(\sum_kp_k=1\).  For the propensity estimate, condition on
\(N_{k,-s}\).  On a nonempty cell its unprojected empirical proportion has
conditional variance at most \(1/(4N_{k,-s})\); on an empty cell its squared
error is at most one.  Projection on
\([\epsilon/2,1-\epsilon/2]\) cannot increase distance from
\(\pi_k\in[\epsilon,1-\epsilon]\).  Repeating (3) gives
\[
 E\!\left[\sum_{k=1}^{d_n}p_k
  (\widehat\pi_k^{(-s)}-\pi_k)^2\right]
 \le C\left\{e^{-c m_s/d_n}+\frac{d_n}{m_s}\right\}.
\tag{5}
\]
Because \(n/d_n\to\infty\), the exponential terms are smaller than any
inverse power of \(n\).  Markov's inequality applied to (4)--(5), for each of
the fixed number of folds, yields
\[
 \|\widehat\mu_a^{(-s)}-\mu_a\|_{P_n,2}
   =O_{Q_{P_n}^{(n)}}(M\sqrt{d_n/n}),
 \qquad
 \|\widehat\pi^{(-s)}-\pi\|_{P_n,2}
   =O_{Q_{P_n}^{(n)}}(\sqrt{d_n/n}).
\tag{6}
\]

Define the uncentered augmented score
\[
 g(O;\eta)=m_1(X)-m_0(X)
 +\frac{A\{Y-m_1(X)\}}{e(X)}
 -\frac{(1-A)\{Y-m_0(X)\}}{1-e(X)},
\tag{7}
\]
where \(\eta=(e,m_0,m_1)\), and let
\(\eta_0=(\pi,\mu_0,\mu_1)\).  The definitions of the conditional means,
fixed overlap, and \(def:ate\text{-}functional\) give
\[
 E_{P_n}[g(O;\eta_0)]=\tau(P_n),
 \qquad g(O;\eta_0)-\tau(P_n)=\phi_{P_n}(O).
\tag{8}
\]
For the nuisance estimate trained outside fold \(s\), put
\(D_s(O)=g(O;\widehat\eta^{(-s)})-g(O;\eta_0)\).  Conditioning on the
training sample and evaluating (7) within each cell gives the exact
orthogonality remainder
\[
 P_nD_s=\sum_{k=1}^{d_n}p_k
 (\widehat\pi_k^{(-s)}-\pi_k)
 \left\{
 \frac{\widehat\mu_{1k}^{(-s)}-\mu_{1k}}
      {\widehat\pi_k^{(-s)}}
 +\frac{\widehat\mu_{0k}^{(-s)}-\mu_{0k}}
      {1-\widehat\pi_k^{(-s)}}
 \right\}.
\tag{9}
\]
Every denominator in (9) is at least \(\epsilon/2\).  Cauchy--Schwarz and
(6) therefore show
\[
 |P_nD_s|=O_{Q_{P_n}^{(n)}}(Md_n/n)
          =o_{Q_{P_n}^{(n)}}(Mn^{-1/2}),
\tag{10}
\]
where the last equality uses \(ass:local\text{-}dimension\).

The same score algebra, \(|\widehat\pi-\pi|\le1\), the clipped denominators,
and (2) imply the conditional second-moment bound
\[
 P_nD_s^2\le C_\epsilon\left[
 \sum_{a=0}^1\sum_{k=1}^{d_n}p_k
 (\widehat\mu_{ak}^{(-s)}-\mu_{ak})^2
 +M^2\sum_{k=1}^{d_n}p_k
 (\widehat\pi_k^{(-s)}-\pi_k)^2\right]
 =O_{Q_{P_n}^{(n)}}(M^2d_n/n).
\tag{11}
\]
Conditionally on the training data, observations in \(I_s\) are independent
of \(\widehat\eta^{(-s)}\).  Hence
\[
 (\mathbb P_{I_s}-P_n)D_s
 =O_{Q_{P_n}^{(n)}}(M\sqrt{d_n}/n).
\tag{12}
\]
Summing (9)--(12) over the fixed number of folds, with weights
\(|I_s|/n\), and using the definition of \(\widehat\tau^{\mathrm{OS}}\), gives
\[
\begin{aligned}
 \widehat\tau^{\mathrm{OS}}-\tau(P_n)
 &=\frac1n\sum_{i=1}^n\{g(O_i;\eta_0)-\tau(P_n)\}
   +O_P(Md_n/n)+O_P(M\sqrt{d_n}/n)\\
 &=\frac1n\sum_{i=1}^n\phi_{P_n}(O_i)
   +o_P(n^{-1/2}).
\end{aligned}
\tag{13}
\]
Here and below the final little-oh is under \(Q_{P_n}^{(n)}\); \(M\) is the
fixed scale indexing the class.

It remains to verify the variance statement.  From
\(|\mu_{ak}|\le M\), \(|\tau(P_n)|\le2M\), fixed overlap, and
\(ass:conditional\text{-}fourth\text{-}moment\), the inequality
\(|x_1+x_2+x_3|^4\le27\sum_{j=1}^3|x_j|^4\) applied to
\(def:efficient\text{-}influence\text{-}function\) yields
\[
 E_{P_n}|\phi_{P_n}(O)|^4\le C_\epsilon M^4.
\tag{14}
\]
Consequently Chebyshev's inequality, (1), and (14) give
\[
 \frac1n\sum_{i=1}^n\phi_{P_n}(O_i)^2
   -V_{\mathrm{eff}}(P_n)=o_P(M^2),
 \qquad
 \overline\phi_n^2=o_P(M^2).
\tag{15}
\]
Furthermore, conditioning fold by fold in (11) and applying Markov's
inequality gives
\[
 \frac1n\sum_{s=0}^{K_0-1}\sum_{i\in I_s}D_s(O_i)^2
 =O_P(M^2d_n/n)=o_P(M^2).
\tag{16}
\]
Writing the fitted score as
\(\widehat\psi_i=\tau(P_n)+\phi_{P_n}(O_i)+D_s(O_i)\) for
\(i\in I_s\), the elementary inequality
\[
 \left|\frac1n\sum_i(u_i+v_i-\overline u-\overline v)^2
       -\frac1n\sum_i(u_i-\overline u)^2\right|
 \le2\left\{\frac1n\sum_i(u_i-\overline u)^2
             \frac1n\sum_i(v_i-\overline v)^2\right\}^{1/2}
      +\frac1n\sum_i(v_i-\overline v)^2
\tag{17}
\]
with \(u_i=\phi_{P_n}(O_i)\), \(v_i=D_s(O_i)\), together with
(15)--(16), proves
\[
 \widehat V^{\mathrm{OS}}-V_{\mathrm{eff}}(P_n)=o_P(M^2).
\tag{18}
\]
Division in (18) is legitimate by the lower bound in (1), and proves
\(\widehat V^{\mathrm{OS}}/V_{\mathrm{eff}}(P_n)\to1\).

Finally, (1) and (14) imply
\[
 \sup_n\frac{E_{P_n}|\phi_{P_n}(O)|^4}
                 {V_{\mathrm{eff}}(P_n)^2}
 \le C_\epsilon/v_0^2<\infty.
\tag{19}
\]
The mean-zero identity required by
\(lem:order\text{-}one\text{-}self\text{-}normalized\text{-}clt\) is (8),
and iid sampling and positive variance are supplied by
\(ass:iid\text{-}sampling\) and (1).  That cited order-one result and (19)
therefore give
\[
 \frac{\sum_{i=1}^n\phi_{P_n}(O_i)}
 {\left\{\sum_{i=1}^n[\phi_{P_n}(O_i)-\overline\phi_n]^2\right\}^{1/2}}
 \Rightarrow N(0,1).
\tag{20}
\]
By (15), the denominator in (20), divided by
\(\sqrt{nV_{\mathrm{eff}}(P_n)}\), converges in probability to one.  For any
\(x\in\mathbb R\) and \(0<\delta<1\), intersecting with the event that this
ratio lies in \([1-\delta,1+\delta]\) sandwiches the distribution function of
the oracle standardized mean between the distribution functions in (20) at
the smaller and larger of \(x/(1-\delta)\) and \(x/(1+\delta)\), up to a
probability tending to zero.
First let \(n\to\infty\) and then \(\delta\downarrow0\); continuity of the
standard-normal distribution function proves
\[
 \frac{1}{\sqrt{nV_{\mathrm{eff}}(P_n)}}
 \sum_{i=1}^n\phi_{P_n}(O_i)\Rightarrow N(0,1).
\tag{21}
\]
Combining (13), (21), and the lower variance bound in (1) proves the final
displayed normal limit and completes the proof of the narrowed theorem.

%%% FIELD partial-rare-tail
The established node \(lem:four\text{-}coordinate\text{-}bias\) supplies the bias
\(C_\epsilon Md/[n\log(en)]\).  The stipulated score, however, does not yet
turn this fact into the requested observable tail radius.  Indeed, for
\[
 \psi(x)=\frac12\log\frac{1+x+x^2/2}{1-x+x^2/2}
\]
one has
\[
 \psi'(x)=\frac{2(2-x^2)}{x^4+4}.
\]
Thus the score is increasing only on \((-\sqrt2,\sqrt2)\) and is decreasing
outside that interval.  More concretely, for any odd block count \(B\ge3\)
and any \(\lambda>0\), consider algebraic block outputs
\(Z_1=\cdots=Z_{B-1}=0\) and \(Z_B=R/\lambda\).  The proposed score is
\[
 S_R(t)=-(B-1)\psi(\lambda t)+\psi(R-\lambda t).
\]
For sufficiently large \(R\), continuity and
\(\psi(x)\downarrow0\) as \(x\to\infty\) give
\[
 S_R(0)>0,\quad S_R(\sqrt2/\lambda)<0,\quad
 S_R((R-\sqrt2)/\lambda)>0,\quad S_R(R/\lambda)<0.
\]
Hence the zero equation has at least three roots.  A retained-range argument
could still repair uniqueness, but no such range or a positive slope modulus
is presently part of the construction, and the unbounded outcome model does
not supply one automatically.

%%% FIELD partial-global-upper
The proved node \(lem:four\text{-}coordinate\text{-}bias\) supplies the
deterministic light-branch bias allowance
\(C_\epsilon Md/[n\log(en)]\).  If measurable light and heavy radii
\(R_{\mathrm L},R_{\mathrm H}\) obeyed, on one common pilot-good event,
\[
 Q_P\{|T_{\mathrm L}-E(T_{\mathrm L}\mid I_0)|>R_{\mathrm L}\mid I_0\}
 \le\alpha/4,
 \qquad
 Q_P\{|T_{\mathrm H}-E(T_{\mathrm H}\mid I_0)|>R_{\mathrm H}\mid I_0\}
 \le\alpha/2,
\]
and had expected radii of orders
\(M\{n^{-1/2}+d/[n\log(en)]\}\) and \(Mn^{-1/2}\), respectively, while
pilot failure and the pilot-dependent conditional means were controlled at
level \(\alpha/4\), then the triangle inequality and a union bound would give
coverage \(1-\alpha\); clipping to \([-2M,2M]\) would not enlarge length, so
integration of the two score tails would give the claimed expected-length
order.  The current core proves only the displayed deterministic bias term.

%%% FIELD partial-finite-precision
Suppose an exact-real interval \([L,U]\) has already been defined and interval
arithmetic returns \(\widetilde L\le L\le U\le\widetilde U\) with
\(L-\widetilde L\le\delta_{\mathrm{num}}\) and
\(\widetilde U-U\le\delta_{\mathrm{num}}\).  Then deterministically
\([L,U]\subseteq[\widetilde L,\widetilde U]\), coverage is monotone under
this inclusion, and
\[
 (\widetilde U-\widetilde L)-(U-L)\le2\delta_{\mathrm{num}}.
\]
The banked coefficient bound \(|g_{Kj}|\le6^{2K}\) with \(K\le L_n\) also
places polynomial/factorial intermediate magnitudes below the declared
\((en)^{8L_n}\) envelope after the frozen input-precision condition is used.
What is missing is a uniform outward enclosure for the roots: the stipulated
score is not monotone and its derivative can vanish, so no uniform
score-to-root error conversion or bisection count follows from the current
formula.
